\subsection{与支撑集的关系}

设 M 是环 A 上的模。回忆，A 的满足 \(M_{p} \neq 0\) 的素理想 p 所成的集合称为 M 的支撑集，记为 Supp(M)（第二章，第 4 节，第 4 节，定义 5）。

命题 7. 设 A 是环，M 是 A-模。

\begin{enumerate}
\def\labelenumi{(\roman{enumi})}
\item
A 的每个素理想 \(\mathfrak{p}\)，若包含 \(\operatorname{Ass}(\mathbf{M})\) 中的一个元素，则属于 \(\operatorname{Supp}(\mathbf{M})\)。
\item
反之，若 A 是诺特环，则每个理想 \(\mathfrak{p} \in \operatorname{Supp}(\mathbf{M})\) 都包含 \(\operatorname{Ass}(\mathbf{M})\) 中的一个元素。
\end{enumerate}

如果 \(\mathfrak{p}\) 包含 \(\mathfrak{q}\) 中的一个元素 \(\operatorname{Ass}(\mathbf{M})\)，则 \(\mathfrak{q} \cap (\mathbf{A} - \mathfrak{p}) = \varnothing\)，因此若记 \(\mathbf{S} = \mathbf{A} - \mathfrak{p}\)，则 \(\mathbf{S}^{-1}\mathfrak{p}\) 是与 \(\mathbf{S}^{-1}\mathbf{M} = \mathbf{M}_{\mathfrak{p}}\) 相伴的素理想（第 2 节，命题 5），特别地 \(\mathbf{M}_{\mathfrak{p}} \neq 0\)，从而 \(\mathfrak{p} \in \operatorname{Supp}(\mathbf{M})\)。反之，若 \(\mathbf{A}\) 是诺特环，则 \(\mathbf{A}_{\mathfrak{p}}\) 也是诺特环（第二章，第 2 节，第 4 节，命题 10 的推论 2）。若 \(\mathbf{M}_{\mathfrak{p}} \neq 0\)，则 \(\operatorname{Ass}_{\mathbf{A}_{\mathfrak{p}}}(\mathbf{M}_{\mathfrak{p}}) \neq 0\)（第 1 节，命题 2 的推论 1），所以存在 \(\mathfrak{q} \in \operatorname{Ass}_{\mathbf{A}}(\mathbf{M})\)，使得 \(\mathfrak{q} \cap (\mathbf{A} - \mathfrak{p}) = \varnothing\)（第 2 节，命题 5 的推论）。

推论 1. 如果 M 是诺特环上的模，则 \(\operatorname{Ass}(\mathbf{M}) \subset \operatorname{Supp}(\mathbf{M})\)，且这两个集合具有相同的极小元。

推论 2. 诺特环 A 的幂零根是理想 \(\mathfrak{p} \in \operatorname{Ass}(\mathbf{A})\) 的交。

我们知道，A 的幂零根是 \(\operatorname{Spec}(A) = \operatorname{Supp}(A)\) 的极小元的交（第二章，第 2 节，第 6 节，命题 13）。

